\documentclass[12pt,a4paper,openany,oneside]{book}

% --- Paquetes ---
\usepackage[utf8]{inputenc}
\usepackage[spanish, es-noshorthands]{babel}
\usepackage{amsmath, amsfonts, amssymb}
\usepackage{graphicx}
\usepackage{geometry}
\usepackage{hyperref}
\usepackage{booktabs}
\usepackage{fancyhdr}
\usepackage{listings}
\usepackage{xcolor}
\usepackage{tikz}
\usepackage{pgfplots}
\usepackage{colortbl}
\usepackage{multirow}
\usepackage{hhline}
\usetikzlibrary{arrows.meta, positioning, shapes.geometric, calc, shapes}
\pgfplotsset{compat=1.18}

\geometry{margin=2.5cm}

% --- Colores y estilos para código Python ---
\definecolor{codegreen}{rgb}{0,0.6,0}
\definecolor{codegray}{rgb}{0.5,0.5,0.5}
\definecolor{codepurple}{rgb}{0.58,0,0.82}
\definecolor{backcolour}{rgb}{0.95,0.95,0.92}

\lstset{
    language=Python,
    backgroundcolor=\color{backcolour},   
    commentstyle=\color{codegreen},
    keywordstyle=\color{blue},
    numberstyle=\tiny\color{codegray},
    stringstyle=\color{codepurple},
    basicstyle=\ttfamily\small,
    breaklines=true,
    frame=single,
    showstringspaces=false
}

% --- Estilo de cabecera y pie ---
\pagestyle{fancy}
\fancyhf{}
\renewcommand{\headrulewidth}{0.4pt}
\renewcommand{\footrulewidth}{0.4pt}
\fancyhead[L]{\small Carlos César Sánchez Coronel}
\fancyhead[R]{\small Sesión 13: Interpretabilidad de Modelos}
\fancyfoot[L]{\small Carlos César Sánchez Coronel}
\fancyfoot[C]{\thepage}

% --- Portada ---
\title{
    \textbf{Sesión 13} \\ 
    \Huge \textbf{INTERPRETABILIDAD DE MODELOS} \\
    \Large Explicando las decisiones de la caja negra
}
\author{
    \small Por \\ \vspace{0.2cm}
    \Large \textsc{Carlos César Sánchez Coronel}
}
\date{2026}

\begin{document}

\maketitle
\tableofcontents
\addcontentsline{toc}{chapter}{Índice general}

% ------------------------------------------------------------------
% CAPÍTULO 13: SESIÓN 13 - INTERPRETABILIDAD DE MODELOS
% ------------------------------------------------------------------
\chapter{Interpretabilidad de Modelos}

En las sesiones anteriores se han construido modelos cada vez más complejos: Random Forest, XGBoost, redes neuronales. Estos modelos logran una alta precisión, pero a costa de convertirse en ``cajas negras'': es difícil entender por qué toman una decisión determinada. En muchos ámbitos (salud, finanzas, justicia), esta falta de transparencia es inaceptable. La interpretabilidad de modelos busca abrir esa caja negra, proporcionando explicaciones comprensibles para los seres humanos. Esta sesión presenta las técnicas más avanzadas para explicar modelos, tanto a nivel global como local, con fundamentos matemáticos y ejemplos prácticos.

\section{Logro de la sesión}
Explicar las predicciones de modelos complejos utilizando técnicas de interpretabilidad global y local, comprendiendo sus fundamentos matemáticos, sus aplicaciones en la industria y sus limitaciones.

\section{Introducción: la necesidad de explicaciones}

Los modelos de machine learning modernos, como XGBoost, redes neuronales o Random Forest, logran una precisión asombrosa, pero su funcionamiento interno es opaco. En aplicaciones críticas, como diagnóstico médico, aprobación de créditos o decisiones judiciales, no basta con una predicción: es necesario justificarla. Además, regulaciones como el GDPR europeo otorgan a los ciudadanos el ``derecho a una explicación'' sobre las decisiones automatizadas que les afectan.

La interpretabilidad también es útil para:
\begin{itemize}
    \item \textbf{Depurar modelos}: identificar sesgos, errores o comportamientos indeseados.
    \item \textbf{Generar confianza}: usuarios y stakeholders aceptan mejor un modelo si entienden sus decisiones.
    \item \textbf{Descubrir conocimiento}: en ciencia, las explicaciones pueden revelar relaciones ocultas en los datos.
\end{itemize}

Existen dos grandes enfoques:
\begin{itemize}
    \item \textbf{Interpretabilidad global}: entender el comportamiento promedio del modelo (qué variables son más importantes, cómo afectan en general).
    \item \textbf{Interpretabilidad local}: explicar una predicción individual (por qué este cliente fue rechazado).
\end{itemize}

Además, los métodos pueden ser:
\begin{itemize}
    \item \textbf{Intrínsecos}: el modelo es interpretable por sí mismo (regresión lineal, árboles pequeños).
    \item \textbf{Post-hoc}: se aplican después de entrenar el modelo (SHAP, LIME).
\end{itemize}

\section{Importancia de características en modelos basados en árboles}

\subsection{Importancia por impureza (MDI)}
En árboles de decisión y Random Forest, la importancia de una característica se calcula sumando las reducciones de impureza (Gini o entropía) que produce en todas las divisiones donde aparece, ponderadas por la proporción de muestras que pasan por ese nodo. Formalmente:

\[
\text{Importancia}(j) = \frac{\sum_{t \in T} \Delta i(t) \cdot n(t)}{\sum_{t \in T} n(t)}
\]

donde \( \Delta i(t) \) es la reducción de impureza en el nodo \(t\), \(n(t)\) el número de muestras en ese nodo, y \(T\) el conjunto de nodos donde se usa la característica \(j\).

\textbf{Limitaciones}:
\begin{itemize}
    \item Puede favorecer variables numéricas con muchas categorías.
    \item No indica dirección del efecto (positivo/negativo).
    \item Puede ser engañosa si hay correlaciones.
\end{itemize}

\subsection{Importancia por permutación}
Idea: si una variable es importante, permutar aleatoriamente sus valores rompe la relación con la variable objetivo, aumentando el error. Se calcula:

\[
\text{Importancia}(j) = \frac{1}{K} \sum_{k=1}^K (\text{error}_{\text{perm}} - \text{error}_{\text{original}})
\]

Es más fiable que MDI y funciona con cualquier modelo. Es el método recomendado en scikit-learn (`permutation_importance`).

\textbf{Ejemplo de código}:
\begin{lstlisting}[language=Python]
from sklearn.inspection import permutation_importance
from sklearn.ensemble import RandomForestClassifier

model = RandomForestClassifier().fit(X_train, y_train)
result = permutation_importance(model, X_test, y_test, n_repeats=10)
sorted_idx = result.importances_mean.argsort()
\end{lstlisting}

\section{Partial Dependence Plots (PDP) e Individual Conditional Expectation (ICE)}

\subsection{Definición matemática}
El PDP muestra cómo cambia la predicción promedio al variar una o dos características, marginalizando sobre el resto. Para una característica \(x_s\):

\[
\text{PDP}(x_s) = \mathbb{E}_{x_c} [f(x_s, x_c)] \approx \frac{1}{n} \sum_{i=1}^n f(x_s, x_c^{(i)})
\]

donde \(x_c\) son las otras variables.

El ICE individualiza este concepto: para cada observación \(i\), se mantienen fijas sus otras variables y se varía \(x_s\), obteniendo una curva por instancia.

\subsection{Interpretación}
\begin{itemize}
    \item PDP suave: relación monótona.
    \item PDP plano: variable poco importante.
    \item PDP con forma compleja: interacciones.
    \item ICE: permite ver heterogeneidad (si todas las curvas son paralelas, no hay interacción; si se cruzan, sí).
\end{itemize}

\subsection{Limitaciones}
\begin{itemize}
    \item Supone independencia entre \(x_s\) y \(x_c\) (puede promediar regiones poco realistas si hay correlación).
    \item Costoso computacionalmente si hay muchas variables.
\end{itemize}

\subsection{Ejemplo de código}
\begin{lstlisting}[language=Python]
from sklearn.inspection import PartialDependenceDisplay

PartialDependenceDisplay.from_estimator(model, X_train, ['duration', 'amount'],
                                        kind='average', grid_resolution=20)
# Para ICE usar kind='individual'
\end{lstlisting}

\section{SHAP (SHapley Additive exPlanations)}

SHAP se basa en los valores de Shapley de la teoría de juegos cooperativos, distribuidos de manera justa entre las características.

\subsection{Valores de Shapley}
En un juego cooperativo con \(p\) jugadores, el valor de Shapley para el jugador \(j\) es:

\[
\phi_j = \sum_{S \subseteq P \setminus \{j\}} \frac{|S|! (p - |S| - 1)!}{p!} [v(S \cup \{j\}) - v(S)]
\]

Aplicado a machine learning:
\begin{itemize}
    \item Los ``jugadores'' son las características.
    \item \(v(S)\) es la predicción promedio condicionada a que las características en \(S\) tomen los valores de la instancia actual (o marginalizadas).
    \item El resultado \(\phi_j\) es la contribución de la característica \(j\) a la diferencia entre la predicción y el valor base (esperado).
\end{itemize}

\textbf{Propiedades deseables}:
\begin{itemize}
    \item \textbf{Eficiencia}: la suma de los \(\phi_j\) más el valor base es igual a la predicción.
    \item \textbf{Simetría}: dos características que contribuyen igual tienen el mismo valor.
    \item \textbf{Linealidad}: si el juego es suma de dos, los valores se suman.
    \item \textbf{Dummy}: si una característica no aporta, su valor es 0.
\end{itemize}

\subsection{SHAP en la práctica}
Calcular exactamente los valores de Shapley es exponencial, por lo que existen aproximaciones:
\begin{itemize}
    \item \textbf{Kernel SHAP}: modelo agnóstico, combina LIME con propiedades Shapley.
    \item \textbf{Tree SHAP}: específico para árboles, mucho más rápido (utiliza los caminos de decisión).
    \item \textbf{Deep SHAP}: para redes neuronales.
\end{itemize}

\subsection{Visualizaciones SHAP}
\begin{itemize}
    \item \textbf{Summary plot}: diagrama de dispersión con todos los valores SHAP por variable. Muestra impacto, dirección y distribución.
    \item \textbf{Dependence plot}: similar a PDP pero con valores SHAP, y puede colorear por interacción.
    \item \textbf{Force plot}: representación local de las contribuciones para una instancia.
    \item \textbf{Waterfall plot}: similar, muestra cómo se construye la predicción desde el valor base.
\end{itemize}

\subsection{Interpretación local vs global}
\begin{itemize}
    \item \textbf{Local}: force plot para un cliente.
    \item \textbf{Global}: summary plot (importancia global + signo) y mean(|SHAP|) como medida de importancia.
\end{itemize}

\subsection{Ejemplo de código con TreeExplainer}
\begin{lstlisting}[language=Python]
import shap

explainer = shap.TreeExplainer(model)
shap_values = explainer.shap_values(X_test)

# Summary plot
shap.summary_plot(shap_values, X_test, feature_names=feature_names)

# Force plot para un individuo
shap.force_plot(explainer.expected_value, shap_values[0,:], X_test.iloc[0,:])
\end{lstlisting}

\section{LIME (Local Interpretable Model-agnostic Explanations)}

LIME explica predicciones individuales aproximando el modelo global con un modelo interpretable (lineal, árbol pequeño) en la vecindad de la instancia.

\subsection{Idea fundamental}
Para una instancia \(x\) a explicar:
\begin{enumerate}
    \item Generar muestras aleatorias alrededor de \(x\) (perturbando las características).
    \item Obtener predicciones del modelo original para esas muestras.
    \item Ponderar las muestras según su proximidad a \(x\) (kernel exponencial).
    \item Ajustar un modelo interpretable (por ejemplo, regresión lineal con regularización Lasso) en ese vecindario.
    \item Los coeficientes del modelo lineal son la explicación.
\end{enumerate}

\subsection{Función de proximidad}
Se usa un kernel exponencial:

\[
\pi_x(z) = \exp\left( -\frac{d(x,z)^2}{\sigma^2} \right)
\]

donde \(d\) puede ser distancia euclídea, coseno, etc.

\subsection{Ventajas y limitaciones}
\begin{itemize}
    \item \textbf{Ventajas}: agnóstico al modelo, rápido, intuitivo.
    \item \textbf{Limitaciones}: inestabilidad (la explicación puede variar con el muestreo), dificultad para definir el tamaño de la vecindad.
\end{itemize}

\subsection{Ejemplo de código}
\begin{lstlisting}[language=Python]
from lime.lime_tabular import LimeTabularExplainer

explainer = LimeTabularExplainer(X_train.values, feature_names=feature_names,
                                 class_names=['bueno', 'malo'], mode='classification')
exp = explainer.explain_instance(X_test.iloc[0].values, model.predict_proba)
exp.show_in_notebook()
\end{lstlisting}

\section{Comparación de métodos de interpretabilidad}

\begin{table}[h]
\centering
\begin{tabular}{|l|c|c|c|c|}
\hline
\textbf{Método} & \textbf{Ámbito} & \textbf{Modelo-agnóstico} & \textbf{Velocidad} & \textbf{Interpretación} \\
\hline
Importancia (permutación) & Global & Sí & Rápida & Ranking de variables \\
PDP & Global & Sí & Media & Relación parcial \\
ICE & Local/Global & Sí & Media & Heterogeneidad \\
SHAP & Local/Global & No (Tree) / Sí (Kernel) & Rápida (Tree) / Lenta (Kernel) & Contribuciones aditivas \\
LIME & Local & Sí & Media & Modelo local \\
\hline
\end{tabular}
\caption{Comparación de métodos de interpretabilidad.}
\end{table}

\textbf{¿Cuándo usar cada uno?}
\begin{itemize}
    \item \textbf{Importancia por permutación}: para un ranking rápido y fiable.
    \item \textbf{PDP/ICE}: para entender la relación funcional entre una variable y la predicción.
    \item \textbf{SHAP}: para explicaciones locales consistentes y un resumen global, especialmente con TreeSHAP (es el estándar de facto).
    \item \textbf{LIME}: cuando se necesita una explicación rápida y no se tiene acceso a SHAP, o para modelos no basados en árboles.
\end{itemize}

\section{Aplicaciones en la industria}

\subsection{Finanzas: scoring crediticio}
\textbf{Problema}: un banco niega un crédito a un cliente. El cliente puede exigir una explicación (GDPR). \\
\textbf{Solución}: usar SHAP para mostrar qué variables contribuyeron negativamente (ingresos bajos, muchas consultas recientes). El force plot se puede incluir en la carta de rechazo.

\subsection{Salud: diagnóstico asistido}
\textbf{Problema}: un modelo predice que un paciente tiene alto riesgo de diabetes. \\
\textbf{Solución}: un PDP muestra que el riesgo aumenta con el IMC y la glucosa, ayudando al médico a entender la recomendación.

\subsection{Recursos humanos: selección de candidatos}
\textbf{Problema}: detectar sesgos del modelo (por ejemplo, si el género es una variable importante). \\
\textbf{Solución}: análisis de importancia SHAP por grupo demográfico para verificar equidad.

\section{Caso integrado: Explicación de un modelo de riesgo crediticio}

\subsection{Datos}
Se utiliza el dataset \textbf{German Credit} (UCI), con 1000 solicitudes, 20 variables (mixtas) y objetivo binario (bueno/malo).

\subsection{Modelo}
Se entrena un XGBoost con validación cruzada, obteniendo un AUC de 0.78.

\subsection{Interpretabilidad global}
\begin{itemize}
    \item \textbf{Importancia por permutación}: las variables más importantes son `duration` (duración del crédito), `amount` (monto) y `age` (edad).
    \item \textbf{PDP}:
        \begin{itemize}
            \item Para `duration`: la probabilidad de mal pagador aumenta con la duración.
            \item Para `amount`: relación positiva hasta cierto punto, luego se estabiliza.
        \end{itemize}
    \item \textbf{SHAP summary plot}:
        \begin{itemize}
            \item Muestra que valores altos de `duration` y `amount` contribuyen a riesgo (SHAP positivo).
            \item La variable `purpose` (propósito) tiene varios niveles con efectos mixtos.
        \end{itemize}
\end{itemize}

\subsection{Interpretabilidad local (un cliente rechazado)}
Supongamos un cliente con:
\begin{itemize}
    \item duration = 48 meses
    \item amount = 5000 €
    \item age = 25 años
    \item purpose = "coche"
    \item other\_installment\_plans = "yes"
\end{itemize}

El modelo predice probabilidad de malo = 0.82 (umbral de rechazo 0.5). El \textbf{SHAP force plot} muestra:
\begin{itemize}
    \item El valor base (predicción media) es 0.35.
    \item `duration=48` aporta +0.25 (por encima de la media).
    \item `amount=5000` aporta +0.15.
    \item `age=25` aporta +0.10 (los jóvenes tienen mayor riesgo según el modelo).
    \item `other\_installment\_plans=yes` aporta +0.07.
    \item El resto de variables aportan -0.05 en total.
    \item Suma: 0.35 + 0.25+0.15+0.10+0.07-0.05 = 0.87 (cercano a 0.82, pequeña diferencia por redondeo).
\end{itemize}

\subsection{Comunicación al cliente}
Se genera un informe en lenguaje natural: "Su solicitud fue rechazada principalmente porque la duración del crédito (48 meses) es superior a la media, lo que incrementa el riesgo. También influyen el monto solicitado (5000 €) y su edad (25 años). Estos factores se combinan con otras variables como tener otros planes de pago."

\subsection{Detección de sesgos}
Se calcula la importancia SHAP promedio por género (variable no incluida en el modelo, pero se puede analizar la distribución de predicciones). Si el modelo discrimina indirectamente, se pueden tomar medidas correctivas.

\section{Limitaciones y buenas prácticas}

\begin{itemize}
    \item \textbf{SHAP es aditivo}, pero las interacciones se reparten entre las variables, pudiendo ser confuso.
    \item \textbf{LIME es inestable}: ejecutar varias veces para verificar consistencia.
    \item \textbf{PDP puede ser engañoso} si las variables están correlacionadas (usar PDP condicionales o ALE plots).
    \item \textbf{No confundir correlación con causalidad}: las explicaciones muestran asociaciones, no causas.
    \item \textbf{Validar con expertos}: la interpretabilidad debe ser validada por conocimiento del dominio.
\end{itemize}

\section{Conexión con el resto del curso}

\begin{itemize}
    \item Los métodos de importancia se basan en métricas vistas en sesiones anteriores.
    \item SHAP utiliza conceptos de teoría de juegos, relacionados con matemáticas discretas.
    \item LIME utiliza regresión lineal (local), vista en sesiones de modelos lineales.
    \item La interpretabilidad es clave para la aceptación de modelos en producción (MLOps, sesión 14).
\end{itemize}

\section{Resumen}

\begin{itemize}
    \item La interpretabilidad puede ser global o local, y los métodos pueden ser intrínsecos o post-hoc.
    \item La importancia de características (permutación) da un ranking global robusto.
    \item PDP e ICE muestran relaciones parciales.
    \item SHAP proporciona explicaciones locales consistentes y un resumen global, siendo el estándar actual.
    \item LIME es una alternativa rápida y agnóstica para explicaciones locales.
    \item Las aplicaciones prácticas incluyen finanzas, salud, detección de sesgos y cumplimiento normativo.
    \item Es fundamental elegir el método según el contexto y validar las explicaciones con expertos.
\end{itemize}

\end{document}